% The MPC slide's figure (spec, slide 3): one real episode on one time axis. Input once in the preamble, after
% vnslides.
%
% \vnmpcfigure{<width>}{<height>} draws, with the plot area <width> by <height>:
%   - the three forecast cones of talk/data/mpc_cone_{1,2,3}.csv as 68/95/99% bands, mean ± 1, 1.96 and 2.576 SD;
%   - the realized value W/W0 from talk/data/mpc_path.csv as the red line;
%   - a coral ring at each re-plan in talk/data/mpc_events.csv, with its label.
% x is the trading day of the episode (day 0 = Oct 10, 2025); the event days are 21, 29 and 50. All files come from
% scripts/30_talk_data.jl.

\definecolor{deepskyblue1}{HTML}{00BFFF}
\definecolor{gray95}{HTML}{F2F2F2}
\definecolor{vnmpcred}{HTML}{FF0000}

% \vnmpcband{<cone>}{<z>}{<opacity>}{<tag>}: the band mean ± z·sd of one cone, filled at that opacity
\newcommand{\vnmpcband}[4]{%
  \addplot[name path=vnmpc#1#4hi, draw=none, forget plot]
    table[col sep=comma, x=day, y expr={\thisrow{mean} + #2*\thisrow{sd}}] {data/mpc_cone_#1.csv};
  \addplot[name path=vnmpc#1#4lo, draw=none, forget plot]
    table[col sep=comma, x=day, y expr={\thisrow{mean} - #2*\thisrow{sd}}] {data/mpc_cone_#1.csv};
  \addplot[fill=deepskyblue1, fill opacity=#3, draw=none, forget plot] fill between[of=vnmpc#1#4hi and vnmpc#1#4lo];
}

% \vnmpccone{<cone>}: the 99, 95 and 68% bands of one cone, widest first so the inner bands layer on top
\newcommand{\vnmpccone}[1]{%
  \vnmpcband{#1}{2.576}{0.13}{a}%
  \vnmpcband{#1}{1.96}{0.17}{b}%
  \vnmpcband{#1}{1.0}{0.24}{c}%
}

% \vnmpcevent{<row>}{<node options>}{<label>}: the ring at the re-plan in row <row> (from 0) of mpc_events.csv, on
% the red line, and its label placed by <node options>. The path has one row per day from day 0, so the value on the
% event day is the path's row of that number. The arrow is set in math: Helvetica Neue has no →.
\newcommand{\vnmpcevent}[3]{%
  \pgfplotstablegetelem{#1}{day}\of\vnmpcEvents
  \let\vnmpcDay\pgfplotsretval
  \pgfplotstablegetelem{\vnmpcDay}{w}\of\vnmpcPath
  \let\vnmpcW\pgfplotsretval
  \edef\vnmpcCmd{%
    \noexpand\node[circle, draw=vnaccent, line width=1.2pt, fill=white, inner sep=0pt, minimum size=5pt]
      at (axis cs:\vnmpcDay,\vnmpcW) {};
    \noexpand\node[#2, font=\noexpand\fontsize{7}{8.4}\noexpand\vnmedium, text=vnslideink, fill=white,
      fill opacity=0.8, text opacity=1, inner sep=1.5pt, rounded corners=1pt]
      at (axis cs:\vnmpcDay,\vnmpcW) {#3};}%
  \vnmpcCmd
}

\newcommand{\vnmpcfigure}[2]{%
  \pgfplotstableread[col sep=comma]{data/mpc_path.csv}\vnmpcPath
  \pgfplotstableread[col sep=comma]{data/mpc_events.csv}\vnmpcEvents
  % the band labels sit at the upper edge of each band where the last cone ends
  \pgfplotstableread[col sep=comma]{data/mpc_cone_3.csv}\vnmpcLast
  \pgfplotstablegetrowsof{\vnmpcLast}%
  \pgfmathtruncatemacro\vnmpcLastRow{\pgfplotsretval - 1}%
  \pgfplotstablegetelem{\vnmpcLastRow}{day}\of\vnmpcLast
  \let\vnmpcEndDay\pgfplotsretval
  \pgfplotstablegetelem{\vnmpcLastRow}{mean}\of\vnmpcLast
  \let\vnmpcEndMean\pgfplotsretval
  \pgfplotstablegetelem{\vnmpcLastRow}{sd}\of\vnmpcLast
  \let\vnmpcEndSd\pgfplotsretval
  % how far below its forecast the band exit was, in SDs, from the band_exit row (row 1; the data script checks the
  % order of the re-plans): the exit is too close to the band edge to see, so the label states it
  \pgfplotstablegetelem{1}{z}\of\vnmpcEvents
  \pgfmathparse{abs(\pgfplotsretval)}%
  \pgfmathprintnumberto[fixed, precision=1, assume math mode=true]{\pgfmathresult}{\vnmpcExitZ}%
  \begin{tikzpicture}
    \begin{axis}[
      width=#1, height=#2, scale only axis,
      clip=false,
      xmin=0, xmax=50, ymin=0.90, ymax=1.25,
      axis background/.style={fill=gray95},
      axis line style={draw=none},
      grid=major, major grid style={draw=white, line width=0.8pt},
      tick style={draw=none},
      xtick={0,21,29,50}, xticklabels={Oct 10, Nov 10, Nov 20, Dec 22},
      ytick distance=0.05,
      tick label style={font=\fontsize{8}{9.6}\selectfont, color=vndottick},
      yticklabel style={/pgf/number format/assume math mode=true, /pgf/number format/fixed,
        /pgf/number format/fixed zerofill, /pgf/number format/precision=2},
      xlabel={2025, trading days}, ylabel={$W/W_0$},
      label style={font=\fontsize{8.5}{10}\selectfont, color=vnslidebody},
      xlabel style={yshift=3pt},
    ]
      % forecast cones, each from the value at its re-plan
      \vnmpccone{1}%
      \vnmpccone{2}%
      \vnmpccone{3}%

      % band labels right of the last cone's end
      \pgfmathsetmacro\vnmpcYa{\vnmpcEndMean + 2.576*\vnmpcEndSd}%
      \pgfmathsetmacro\vnmpcYb{\vnmpcEndMean + 1.96*\vnmpcEndSd}%
      \pgfmathsetmacro\vnmpcYc{\vnmpcEndMean + 1.0*\vnmpcEndSd}%
      \node[anchor=west, xshift=3pt, font=\fontsize{7}{8.4}\selectfont, text=vnslidebody]
        at (axis cs:\vnmpcEndDay,\vnmpcYa) {99\%};
      \node[anchor=west, xshift=3pt, font=\fontsize{7}{8.4}\selectfont, text=vnslidebody]
        at (axis cs:\vnmpcEndDay,\vnmpcYb) {95\%};
      \node[anchor=west, xshift=3pt, font=\fontsize{7}{8.4}\selectfont, text=vnslidebody]
        at (axis cs:\vnmpcEndDay,\vnmpcYc) {68\%};

      % the realized value
      \addplot[color=vnmpcred, line width=1.8pt, line join=round, mark=none]
        table[col sep=comma, x=day, y=w] {data/mpc_path.csv};

      % the re-plans
      \vnmpcevent{0}{anchor=south, yshift=7pt}{Nov 10: 21 days up $\rightarrow$ re-plan}%
      \vnmpcevent{1}{anchor=north, yshift=-7pt}{Nov 20: \vnmpcExitZ{} SD below, past the 95\% band $\rightarrow$ re-plan}%
      \vnmpcevent{2}{anchor=south east, xshift=-3pt, yshift=7pt}{Dec 22: 21 days up $\rightarrow$ re-plan}%
    \end{axis}
  \end{tikzpicture}%
}
